\documentclass[12pt]{article}
% Préambule commun à tous les documents extraits de « La Revue CAPES »
\usepackage[T1]{fontenc}
\usepackage[utf8]{inputenc}
\usepackage{lmodern}
\usepackage{amsmath,amssymb,amsthm,amsfonts,mathrsfs}
\usepackage{setspace}
\usepackage{listings}
\usepackage{pgfplots}
\pgfplotsset{compat=1.18}
\usepackage{longtable}
\usepackage{adjustbox}
\usepackage{lettrine}
\usepackage{tkz-tab}
\usepackage{changepage}
\usepackage{pdflscape}
\usepackage{graphicx}
\usepackage{tikz}
\usepackage{xcolor}
\usepackage[a4paper,top=2cm,bottom=2.5cm,left=2.5cm,right=2cm]{geometry}
\usepackage{fancyhdr}
\usepackage{draftwatermark}
\SetWatermarkText{La RevueCapes}
\SetWatermarkLightness{0.9}
\SetWatermarkScale{2}
\usepackage{hyperref}
\hypersetup{colorlinks=true,linkcolor=blue!50!black,urlcolor=blue!50!black}

\definecolor{revue}{RGB}{20,60,120}

\pagestyle{fancy}
\fancyhf{}
\renewcommand{\headrulewidth}{0.4pt}
\setlength{\headheight}{15pt}

% \entete{Matière}{Titre}{Type de document}
\newcommand{\entete}[3]{%
  \fancyhead[L]{\small\textsc{La Revue CAPES} -- #1}%
  \fancyhead[R]{\small #2}%
  \fancyfoot[C]{\thepage}%
  \thispagestyle{plain}%
  \begin{center}
    {\color{revue}\rule{\textwidth}{1.2pt}}\\[4pt]
    {\small\textsc{Concours directs CAP-CEG / CAPES -- Burkina Faso}}\\[6pt]
    {\LARGE\bfseries\color{revue} #2}\\[6pt]
    {\large\itshape #1 -- #3}\\[2pt]
    {\color{revue}\rule{\textwidth}{1.2pt}}
  \end{center}
  \vspace{0.5cm}%
}

% Encadré pour les remarques ajoutées dans les corrigés
\newcommand{\remarque}[1]{\par\medskip\noindent\fcolorbox{revue}{revue!6}{\parbox{\dimexpr\linewidth-2\fboxsep-2\fboxrule}{\textbf{\color{revue}Remarque.} #1}}\par\medskip}


\begin{document}
\entete{Mathématiques}{Sujet type n°10}{Énoncé}
\begin{spacing}{1.5}

\textbf{Exercice 1}

On considère le système $(S)$ suivant, composé de deux équations à deux inconnues 
réelles $x$ et $y$ :
\medskip
\[\begin{cases}
ax-by=c\\
x-2y=3
\end{cases}\]

 Les paramètres $a, b$ et $c$ sont obtenus en lançant trois fois consécutives, de façon indépendante, un dé supposé parfait, à six faces numérotées de 1 à 6. Le premier lancer donne la valeur de $a$, le deuxième celle de $b$ et le troisième celle de $c$. 

1) Combien y a-t-il de valeurs possibles pour le triplet $(a, b, c)$ ? 

Dans toute la suite du problème, on donnera les probabilités demandées sous 
forme de fractions de dénominateur 108. 

2) Calculer la probabilité $P_1$ que $(S)$ ait une infinité de solutions. 

3) Calculer la probabilité $P_2$ que $(S)$ n'admette aucune solution. 

4) Calculer la probabilité $P_3$ que $(S)$ admette une solution unique. 

5) Calculer la probabilité $P_4$ que $(S)$ admette le couple $(x = 3, y = 0)$ comme unique 
solution. 

\medskip\textbf{Exercice 2 }

On considère la série entière:

\[\sum_{n=1}^{+\infty}\frac{x^n}{n(n+2)}\]

1) Déterminer son rayon de convergence $R$.

2) Étudier la nature de la série pour $x=1$ et $x=-1$.

3) Décomposer en éléments simples la fraction rationnelle $\frac{1}{X(X+2)}$

4) Rappeler le développement en série entière de la fonction $\longmapsto ln(1-x)$ en 0.

5) $\forall x\in ]-R,R[$, en déduire la somme:

$S=\sum_{n=1}^{+\infty}\frac{x^n}{n+2}$.

6) Montrer que $\forall x\in ]-R,0[\cup ]0,R[$, $\sum_{n=1}^{+\infty}\frac{x^n}{n(n+2)}=\frac{2(1-x^2)ln(1-x)+2x+x^2}{4x^2}$.

7) Via un équivalent de $\ln(1-x)$, vérifier que l'on peut prolonger cette dernière expression en 0.

8) En utilisant la question précédentes calculer les sommes suivantes :

\[S_1=\sum_{n=1}^{+\infty}\frac{(-1)^n}{n(n+2)} \text{ et } S_2=\sum_{n=1}^{+\infty}\frac{1}{n(n+2)}\]

\medskip\textbf{Exercice 3 }

Soit $A$ la matrice de $M_3(\mathbb{R})$ suivante :

			$A=\begin{pmatrix}
			1 & 0 & 1 \\ 
			-1 & 2 & 1 \\ 
			1 & -1 & 1
			\end{pmatrix} $
			
1. Démontrer que les valeurs propres de $A$ sont 1 et 2

2. Déterminer les sous-espaces propres de $A$. La matrice $A$ est-elle diagonalisable?

3. Déterminer les sous-espaces caractéristiques de $A$.

4. Déterminer une base de $\mathbb{R}^3$ dans laquelle la matrice de l'endomorphisme associés à $A$ est

$B=\begin{pmatrix}
2 & 0 & 0 \\ 
0 & 1 & 1 \\ 
0 & 0 & 1
\end{pmatrix}$

En déduire la décomposition de Dunford de $B$.

5. Résoudre le système différentiel :
\[
\begin{cases}
x'=x+z\\
y'=-x+2y+z\\
z'=x-y+z
\end{cases}
\]

\medskip\textbf{Exercice 4 }

Trouver toutes les applications $f :\mathbb{R} \longrightarrow \mathbb{R}$ deux fois dérivables telles que: 
$\forall x \in \mathbb{R} , f"(x) + f(-x) = x \hspace{0.1cm}(E)$

\underline{Indication} : On pourra poser
$g(x) = f(x) + f(-x)$ et $h(x) = f(x) -f(-x)$ et établir deux équations différentielles linéaires du second ordre dont $g$ et $h$ sont solutions.


\quad
\quad

\end{spacing}
\end{document}
